We have introduced \emph{permABC}, a novel framework for Approximate Bayesian Computation (ABC) tailored to hierarchical models with exchangeable structure. By exploiting the invariance of compartments under permutations, permABC restores identifiability among local parameters and significantly improves acceptance rates in high-dimensional settings.

Extending this framework, we developed a sequential version, \emph{permABC-SMC}, which adapts classical ABC-SMC algorithms to permutation-based acceptance rules. We further proposed two complementary strategies: \emph{Over-Sampling}, which enriches the simulation space early on to facilitate acceptance, and \emph{Under-Matching}, which relaxes the matching constraint to allow robust early-stage inference. These approaches provide new tools for navigating pseudo-posterior sequences, enabling both computational efficiency and resilience to model misspecification.

Our theoretical results show that permABC retains the desirable convergence properties of classical ABC algorithms in the small-tolerance regime. Furthermore, numerical experiments highlight the efficiency and stability of permABC in high-dimensional scenarios, as well as its robustness to model misspecification where coordinate-wise approaches like ABC-Gibbs may fail.

\new{We also clarified the connection between permABC and Wasserstein ABC: at the level of compartment-level empirical measures, permABC is a natural hierarchical analogue of WABC, and in the one-dimensional case the two approaches coincide. Asymptotic arguments suggest a WABC-style concentration on the global parameter as \(K \to \infty\), while the over-sampling limit \(M \to \infty\) is a qualitatively different object acting as an algorithmic bridge rather than an inferential target. On the computational side, warm-started swap refinement reduces the amortised cost of the assignment step within SMC from \(\mathcal{O}(K^3)\) to \(\mathcal{O}(K^2)\) per particle.}



Finally, we demonstrated the practical relevance of our method through an application to COVID-19 epidemic data, using a hierarchical SIR model across 94 departments. This application highlights the practical impact of the permABC: leveraging local heterogeneity leads to sharper inference on global parameters, opening new directions for scalable and robust likelihood-free inference.

All proposed methods are implemented in a dedicated Python package called \texttt{permabc}, available at \url{https://github.com/AntoineLuciano/permABC}. The repository also includes code to reproduce the figures and experiments presented in this paper.


Future directions include a deeper theoretical and empirical investigation of the Over-Sampling and Under-Matching strategies, with particular emphasis on the design of more efficient proposal distributions. An additional promising avenue is to explore joint or adaptive use of both strategies within a unified sequential framework: their complementarity may further enhance inference performance.
